\documentclass[11pt,a4paper]{article}

% ---------------------------------------------------------
% PACKAGES
% ---------------------------------------------------------
\usepackage[utf8]{inputenc}
\usepackage[T1]{fontenc}
\usepackage{lmodern}
\usepackage{amsmath, amssymb, amsfonts, bm}
\usepackage{geometry}
\usepackage{graphicx}
\usepackage{hyperref}
\usepackage{authblk}
\usepackage{cite}
\usepackage{physics}
\usepackage{mathtools}
\usepackage{tensor}
\usepackage{enumitem}

\geometry{margin=2.4cm}

% ---------------------------------------------------------
% TITLE
% ---------------------------------------------------------
\title{\textbf{Companion IX (v49.1)}\\[0.2cm]
\textbf{Parameter Estimation and RDCC Parameter Space}\\[0.2cm]
\large Foundations of Relaxation-Driven Cyclic Cosmology (RDCC)}

\author[1]{Michael Lehmann}
\affil[1]{\small Independent Researcher, Relaxation-Driven Cyclic Cosmology (RDCC)\\
\texttt{mi.lehmann@gmx.de}}

\date{September 2026}

\begin{document}
\maketitle

% ========================================================
% ABSTRACT
% ========================================================

\begin{abstract}
This Companion develops the corrected parameter-estimation structure of RDCC in
its Phase-1-integrated v49.1 form. The RDCC parameter space is globally defined
and incorporates the infrared relaxation parameter $\alpha_{\mathrm{IR}}$, the
mirror-sector density contrast $\delta_{(-)}$, the geometric coefficients
$C_1$–$C_4$, and the ultraviolet scale $M$. The likelihood structure developed
in Companion VIII is used to construct the posterior distribution, and the
Boltzmann hierarchy of Companion VII provides the theoretical predictions. This
Companion establishes the statistical backbone of RDCC and defines the
parameter space used in all survey analyses.
\end{abstract}

\tableofcontents
\newpage
% ========================================================
% 1 INTRODUCTION
% ========================================================

\section{Introduction}

Parameter estimation in RDCC differs fundamentally from $\Lambda$CDM. RDCC is a
global framework defined on a bipartite Hilbert space, and its parameters
reflect the geometric, thermodynamic, and ultraviolet structure of the global
state.

In v49.1, the parameter space is corrected to:
\begin{itemize}[leftmargin=0.8cm]
  \item include the infrared relaxation parameter $\alpha_{\mathrm{IR}}$,
  \item include the mirror-sector density contrast $\delta_{(-)}$,
  \item treat $C_1$–$C_4$ as fixed theoretical factors,
  \item incorporate the ultraviolet scale $M$,
  \item remain fully EFT-coherent and statistically consistent.
\end{itemize}

This Companion develops the corrected RDCC parameter space and its connection to
the likelihood structure of Companion VIII.
% ========================================================
% 2 RDCC PARAMETER SPACE
% ========================================================

\section[RDCC Parameter Space]{RDCC Parameter Space}

The RDCC parameter space consists of:
\begin{itemize}[leftmargin=0.8cm]
  \item standard cosmological parameters,
  \item RDCC-specific parameters,
  \item survey nuisance parameters.
\end{itemize}

\subsection[Standard Parameters]{Standard Cosmological Parameters}

\begin{align}
\Omega_{\mathrm{b}} h^2 &\quad \text{baryon density}, \\
\Omega_{\mathrm{r}} h^2 &\quad \text{radiation density}, \\
H_0 &\quad \text{Hubble parameter}, \\
n_s &\quad \text{scalar spectral index}, \\
A_s &\quad \text{scalar amplitude}.
\end{align}

\subsection[RDCC-Specific Parameters]{RDCC-Specific Parameters}

\begin{align}
\alpha_{\mathrm{IR}} &\quad \text{infrared relaxation parameter}, \\
M &\quad \text{ultraviolet suppression scale}, \\
\delta_{(-)} &\quad \text{mirror-sector density contrast}, \\
C_1, C_2, C_3, C_4 &\quad \text{fixed geometric coefficients}.
\end{align}

\subsection[Nuisance Parameters]{Survey Nuisance Parameters}

These include calibration parameters, bias parameters, and noise parameters for
each survey.
% ========================================================
% 3 POSTERIOR DISTRIBUTION
% ========================================================

\section[Posterior Distribution]{Posterior Distribution}

The posterior distribution is defined as
\begin{equation}
P(\boldsymbol{\theta}|\mathbf{d})
\propto
\mathcal{L}(\boldsymbol{\theta})
\,
\pi(\boldsymbol{\theta}),
\end{equation}
where:
\begin{itemize}[leftmargin=0.8cm]
  \item $\mathcal{L}(\boldsymbol{\theta})$ is the likelihood from Companion VIII,
  \item $\pi(\boldsymbol{\theta})$ is the prior distribution.
\end{itemize}

The likelihood is
\begin{equation}
\mathcal{L}(\boldsymbol{\theta})
=
\exp\left(
-\frac{1}{2}
\Delta\mathbf{d}^{\mathrm{T}}
\mathbf{C}^{-1}
\Delta\mathbf{d}
\right),
\end{equation}
with residual vector
\begin{equation}
\Delta\mathbf{d}
=
\mathbf{d}
-
\mathbf{t}(\boldsymbol{\theta}).
\end{equation}
% ========================================================
% 4 PRIORS
% ========================================================

\section{Priors}

The priors reflect the global geometric structure of RDCC.

\subsection[Standard Priors]{Standard Cosmological Priors}

Flat or weakly informative priors are used for:
\begin{align}
\Omega_{\mathrm{b}} h^2, \quad
\Omega_{\mathrm{r}} h^2, \quad
H_0, \quad
n_s, \quad
A_s.
\end{align}

\subsection[RDCC Priors]{RDCC-Specific Priors}

\begin{itemize}[leftmargin=0.8cm]
  \item $\alpha_{\mathrm{IR}}$: positive, small, infrared-saturated,
  \item $M$: large ultraviolet scale,
  \item $\delta_{(-)}$: diffuse, non-compact halo contrast,
  \item $C_1$–$C_4$: fixed theoretical factors of order unity.
\end{itemize}
% ========================================================
% 5 SAMPLING AND INFERENCE
% ========================================================

\section{Sampling and Inference}

Sampling is performed using standard MCMC or nested sampling algorithms. The
RDCC-specific parameters enter through:
\begin{itemize}[leftmargin=0.8cm]
  \item the Boltzmann hierarchy (Companion VII),
  \item the relaxation geometry (Companion VI),
  \item the tensor sector (Companion II and XI),
  \item the thermodynamic sector (Companion IV),
  \item the ultraviolet sector (Companion X).
\end{itemize}

The posterior distribution provides constraints on:
\begin{itemize}[leftmargin=0.8cm]
  \item $\alpha_{\mathrm{IR}}$,
  \item the mirror-sector density contrast $\delta_{(-)}$,
  \item the ultraviolet scale $M$,
  \item the geometric coefficients $C_1$–$C_4$.
\end{itemize}
% ========================================================
% 6 CONCLUSION
% ========================================================

\section{Conclusion}

This Companion develops the corrected parameter-estimation structure of RDCC in
its v49.1 form. The RDCC parameter space incorporates the infrared relaxation
parameter $\alpha_{\mathrm{IR}}$, the mirror-sector density contrast
$\delta_{(-)}$, the geometric coefficients $C_1$–$C_4$, and the ultraviolet
scale $M$. The likelihood structure of Companion VIII and the Boltzmann
hierarchy of Companion VII provide the theoretical backbone for parameter
estimation.

These results integrate the relaxation sector (Companion III and VI), the
thermodynamic sector (Companion IV), the ultraviolet sector (Companion X), and
the survey interface (Companion XII). They form the statistical foundation of
RDCC and connect the theoretical framework to observational data.

\end{document}
